% !TeX program = xelatex
% Вопрос 2. Только цитаты из ../tex/2/chapter.tex.
% Библиографические сведения: ../tex/2/main.tex.
% Источником отбора цитат служит TeX, а не иные книги или сведения по памяти.
\documentclass[12pt,a4paper]{article}
\usepackage{fontspec}
\setmainfont{Times New Roman}
\usepackage[english,russian]{babel}
\usepackage{amsmath,amssymb}
\usepackage[margin=22mm]{geometry}
\usepackage[unicode,hidelinks]{hyperref}
\DeclareMathOperator{\grad}{grad}
\DeclareMathOperator{\Div}{div}
\newcommand{\Rey}{\mathrm{Re}}
\setlength{\emergencystretch}{3em}
\setlength{\parindent}{0pt}
\renewcommand{\thefootnote}{\arabic{footnote})}
\newenvironment{textquote}{\begin{quote}}{\end{quote}\nopagebreak[4]}
\newcommand{\source}[2]{\par\nobreak\smallskip{\small Н. Е. Леонтьев. \textit{Основы теории фильтрации}. 2-е изд., 2017. Глава I, §~#1, с.~#2.}\par\medskip}
\begin{document}
% BEGIN QUESTION HEADER
\section*{Вопрос 2}
\noindent Закон Дарси и пределы его применимости. Учёт анизотропии пористой среды. Двучленный закон фильтрации Форхгеймера. Фильтрация с предельным градиентом давления.

\medskip
\noindent\textbf{Источник:} Н. Е. Леонтьев. \textit{Основы теории фильтрации}. 2-е изд. Москва: МАКС Пресс, 2017. Глава I, §§ 3--4, печатные страницы 13--24. Ответ составлен из дословных цитат перенесённого в TeX текста учебника.
\par\bigskip
% END QUESTION HEADER
% Цитата 1; печатные страницы: 14--15
\begin{textquote}
Исторически уравнение баланса импульса было экспериментально получено французским гидравликом Анр\'и Дарс\'и (H.Darcy, 1856) для случая медленного стационарного движения несжимаемой жидкости в неподвижной изотропной пористой
среде. В современных обозначениях это соотношение, называемое \textit{законом Дарси}, для фильтрации в поле тяжести имеет вид
\begin{equation}
0=-\grad p+\rho\boldsymbol{g}-\frac{\mu}{k}\boldsymbol{u},
\tag{3.3}\label{q2:eq-3-3}
\end{equation}
где $\boldsymbol{g}$ --- ускорение свободного падения, $\mu$ --- коэффициент динамической вязкости жидкости, $k$ --- коэффициент, называемый проницаемостью пористой среды, который зависит от типа пористой среды. По структуре уравнение \eqref{q2:eq-3-3} похоже на уравнение Навье--Стокса для вязкой несжимаемой жидкости при $\mu=\mathrm{const}$:
\begin{equation}
\rho\frac{d\boldsymbol{v}}{dt}=-\grad p+\rho\boldsymbol{g}+\mu\triangle\boldsymbol{v}
\tag{3.4}\label{q2:eq-3-4}
\end{equation}
и это не удивительно: макроскопический закон движения при фильтрации в пористой среде (закон Дарси) может быть получен с помощью некоторой процедуры осреднения из уравнения движения, описывающего внутрипоровое движение жидкости на микроуровне, то есть уравнения Навье--Стокса (см. § 5).
\source{3}{14--15}
\end{textquote}

% Цитата 2; печатные страницы: 15
\begin{textquote}
При медленных (ползущих) течениях внутри пористой среды, для которых и справедлив закон Дарси, инерционные силы много меньше вязких сил, поэтому сто\'ящим в левой части \eqref{q2:eq-3-4} ускорением можно пренебречь, что даст при осреднении ноль в левой части \eqref{q2:eq-3-3}. Наличие же последнего слагаемого в \eqref{q2:eq-3-4}, описывающего действие вязкости, приводит при переходе к макроскопическому описанию к появлению объемной силы вязкого сопротивления, действующей на жидкость со стороны пористого скелета. В рассматриваемом случае медленных течений\footnote[1]{Более точно --- при числах Рейнольдса $\Rey=\rho vd/\mu\ll1$.} в изотропной среде сила сопротивления направлена противоположно скорости фильтрации, а ее величина прямо пропорциональна вязкости и модулю скорости фильтрации.
\source{3}{15}
\end{textquote}

% Цитата 3; печатные страницы: 15
\begin{textquote}
При $\boldsymbol{g}=0$ закон Дарси приобретает вид
\begin{equation}
\boldsymbol{u}=-\frac{k}{\mu}\grad p;
\tag{3.5}\label{q2:eq-3-5}
\end{equation}
знак минус в правой части \eqref{q2:eq-3-5} имеет ясный физический смысл: в отсутствие массовых сил жидкость движется в пористой среде из областей с б\'ольшим давлением в области с меньшим давлением.
\source{3}{15}
\end{textquote}

% Цитата 4; печатные страницы: 16
\begin{textquote}
Коэффициент проницаемости $k$ зависит только от свойств пористой среды (но не от свойств жидкости), и определяется, в основном, геометрией порового пространства.

[…]

Для неоднородных пористых сред коэффициент проницаемости является функцией пространственных координат и может меняться с течением времени (например, при деформациях пористого скелета или за счет его засорения или растворения).
\source{3}{16}
\end{textquote}

% Цитата 5; печатные страницы: 17
\begin{textquote}
При фильтрации в анизотропных пористых средах (например, в осадочных слоистых пластах или в средах с трещинами, направленными в выделенном направлении) линейная зависимость скорости фильтрации и градиента давления сохраняется (для медленных течений), однако эти векторы в случае общего положения уже не будут параллельны. Проницаемость таких сред будет характеризоваться не одним числом (коэффициентом проницаемости $k$), а симметричным положительно-определенным \textit{тензором проницаемости} $k_{ij}$, входящим в закон Дарси
\[
u_i=-\frac{k_{ij}}{\mu}(\nabla^jp-\rho g^j).
\]
Например, в слоистой среде, у которой направление, перпендикулярное плоскости напластования, задается единичным вектором $\boldsymbol{n}$, тензор проницаемости в декартовой системе координат имеет вид
\[
k_{ij}=(k_\perp-k_\parallel)n_in_j+k_\parallel\delta_{ij},
\]
где $k_\perp$ и $k_\parallel$ суть проницаемости при фильтрации соответственно поперек и вдоль напластования.
\source{3}{17}
\end{textquote}

% Цитата 6; печатные страницы: 18
\begin{textquote}
В некоторых практически важных ситуациях при фильтрации жидкостей --- речь идет об однородных жидкостях, целиком заполняющих недеформируемые изотропные пористые среды --- наблюдаются отступления от классического закона Дарси \eqref{q2:eq-3-3}. Эти отклонения связаны с двумя причинами. \textit{Первая} --- это немалость скоростей движения жидкости внутри пор, т.е. невозможность пренебречь инерционными членами в микроскопических уравнениях \eqref{q2:eq-3-4}. Такие условия, в которых неприменимо приближение Стокса, возникают, например, в прискважинных областях или при фильтрации в крупнозернистых пористых средах (фильтрация через гравий и гальку, а также через гранулированные среды в промышленных технологических установках). \textit{Вторая} причина состоит в отличии реологических\footnote[1]{Реология --- наука, изучающая связи между напряжениями и деформационными характеристиками материалов. Название науки, предложенное американским исследователем Ю.Бингамом (E.Bingham, 1920), происходит от греческого $\rho\varepsilon\omega$ --- «теку», «струюсь».} соотношений фильтрующейся жидкости от закона Навье--Стокса, причем неньютоновское поведение жидкостей при течении внутри порового пространства может проявляться как в узкой пристеночной области (за счет физико-химического взаимодействия жидкости со стенкой поры), так и во всем объеме жидкости (например, при фильтрации некоторых нефтей).
\source{4}{18}
\end{textquote}

% Цитата 7; печатные страницы: 18--19
\begin{textquote}
Предполагая для простоты, что массовые силы отсутствуют\footnote[2]{Влияние потенциальных сил можно учесть путем переопределения давления (например, в случае фильтрации несжимаемой однородной жидкости в поле силы тяжести --- с помощью введения новой переменной $p'=p-\rho\boldsymbol{r}\cdot\boldsymbol{g}$, где $\boldsymbol{r}$ --- радиус-вектор).}, установим общий вид зависимости градиента давления от скорости фильтрации с использованием теории размерностей.

Модуль градиента давления $|\grad p|$ при стационарной фильтрации зависит, кроме модуля скорости фильтрации $u$, от свойств жидкости (плотности $\rho$ и вязкости $\mu$) и пористой среды (характерного размера пор $d$, пористости $m$ и, возможно, других скалярных безразмерных параметров $a_\alpha$, определяющих геометрию
порового пространства):
\[
|\grad p|=F(u,\rho,\mu,d,m,a_\alpha),
\]
и применение $\pi$-теоремы [6,15] дает зависимость
\[
|\grad p|=\frac{\mu u}{d^2}\cdot\Phi(\Rey,m,a_\alpha),\qquad \Rey=\frac{\rho ud}{\mu}.
\]
\source{4}{18--19}
\end{textquote}

% Цитата 8; печатные страницы: 19--20
\begin{textquote}
Итак, общий вид закона фильтрации несжимаемой жидкости в изотропной пористой среде можно записать так:
\begin{equation}
\grad p=-\frac{\mu\boldsymbol{u}}{d^2}\Phi(\Rey,m,a_\alpha).
\tag{4.1}\label{q2:eq-4-1}
\end{equation}
Зависимость $\Phi(\Rey)$ в двух предельных случаях --- при малых и больших числах Рейнольдса --- можно установить из общих соображений.

В случае \textit{очень медленных} (ползущих) внутрипоровых течений, в которых доминируют вязкие эффекты, зависимость \eqref{q2:eq-4-1}
не должна содержать плотность жидкости, выпадающую из уравнений вязкой жидкости в приближении Стокса. Поэтому при $\Rey\ll1$ функция $\Phi$ не должна зависеть от числа Рейнольдса, и закон фильтрации \eqref{q2:eq-4-1} можно записать в виде
\[
\grad p=-\frac{\mu}{k}\boldsymbol{u},\quad k=\frac{d^2}{\Phi(0,m,a_\alpha)},
\]
где величина $k$, имеющая размерность квадрата длины и зависящая только от свойств пористой среды, должна, в силу закона Дарси \eqref{q2:eq-3-5}, совпадать с коэффициентом проницаемости.
\source{4}{19--20}
\end{textquote}

% Цитата 9; печатные страницы: 20
\begin{textquote}
В противоположном случае \textit{быстрых внутрипоровых течений} (при $\Rey\gg1$) преобладающими должны быть инерционные эффекты, поэтому зависимость градиента давления от скорости фильтрации не должна содержать в главных членах коэффициент вязкости $\mu$. Это дает при больших $\Rey$ линейную зависимость $\Phi(\Rey)$ и квадратичную зависимость $|\grad p|$ от $u$.

Если нормировать функцию $\Phi(\Rey)$, вводя обозначение
\begin{equation}
f(\Rey)=\Phi(\Rey)/\Phi(0),\quad f(0)=1,\quad f(\Rey)=O(\Rey)\text{ при }\Rey\to\infty,
\tag{4.2}\label{q2:eq-4-2}
\end{equation}
то закон фильтрации в принимает окончательный вид
\[
0=-\grad p+\rho\boldsymbol{g}-\frac{\mu}{k}\boldsymbol{u}\cdot f(\Rey),\quad \Rey=\frac{\rho ud}{\mu}.
\]
\source{4}{20}
\end{textquote}

% Цитата 10; печатные страницы: 21
\begin{textquote}
При практических расчетах закон $f(\Rey)$ часто аппроксимируют линейной зависимостью, удовлетворяющей требованиям \eqref{q2:eq-4-2}:
\[
f(\Rey)=1+\frac{\beta\sqrt{k}}{d}\Rey,
\]
где $\beta$ --- безразмерный параметр, зависящий от структуры пористой среды (порядка единицы). Это дает \textit{двучленный закон фильтрации}, называемый также \textit{законом фильтрации Форхгеймера} (P.Forchheimer, 1901):
\begin{equation}
0=-\grad p+\rho\boldsymbol{g}-\frac{\mu}{k}\boldsymbol{u}-\beta\frac{\rho u^2}{\sqrt{k}}\cdot\frac{\boldsymbol{u}}{u}.
\tag{4.3}\label{q2:eq-4-3}
\end{equation}
\source{4}{21}
\end{textquote}

% Цитата 11; печатные страницы: 21
\begin{textquote}
Как видно из рис.~5, отклонения от закона Дарси, которому соответствует горизонтальный участок графика, начинаются при $\Rey$ порядка единиц (числа Рейнольдса вычисляются по характерному размеру зерен). При дальнейшем повышении скорости фильтрации инерцией жидкости при ее движении в извилистом внутрипоровом пространстве уже нельзя пренебрегать и нужно пользоваться, например, квадратичной зависимостью \eqref{q2:eq-4-3}. Этот закон по своей структуре (квадратичная зависимость от скорости) похож на известную связь сопротивления и скорости при развитом турбулентном движении жидкости в каналах или при обтекании тел потоком. Основываясь на этой формальной аналогии, иногда режим фильтрации в области применимости закона \eqref{q2:eq-4-3} называют «турбулентной фильтрацией». В действительности же, как показывают эксперименты, при переходе через границу применимости закона Дарси движение жидкости в порах остается ламинарным, а переход к турбулентности происходит при достаточно больших числах Рейнольдса (около 100--300).
\source{4}{21}
\end{textquote}

% Цитата 12; печатные страницы: 22--23
\begin{textquote}
В ряде природных процессов, например при просачивании воды в мелкопористых средах или при фильтрации некоторых нефтей, течение жидкости начинается при превышении модулем градиента давления $|\grad p|$ некоторого фиксированного порогового
значения $G$ («предельного градиента»), а при меньших градиентах жидкость в порах покоится\footnote[1]{Отметим для сравнения, что при выполнении закона Дарси фильтрационное течение начинается при любых сколь угодно малых градиентах давления.}.
\source{4}{22--23}
\end{textquote}

% Цитата 13; печатные страницы: 23
\begin{textquote}
Для моделирования таких явлений часто используется \textit{закон фильтрации с предельным (начальным) градиентом}
\begin{equation}
\boldsymbol{u}=\begin{cases}
-\dfrac{k}{\mu}(|\grad p|-G)\cdot\dfrac{\grad p}{|\grad p|}&\text{при }|\grad p|\geqslant G,\\[6pt]
0&\text{при }|\grad p|<G.
\end{cases}
\tag{4.4}\label{q2:eq-4-4}
\end{equation}
Особенностью этой модели является возможность существования застойных зон (с нулевой скоростью фильтрации), в которых давление не определено однозначно и зависит от предыстории течения.
\source{4}{23}
\end{textquote}

% Цитата 14; печатные страницы: 23--24
\begin{textquote}
Наличие предельного градиента легко объяснить, если предположить, что фильтрующаяся жидкость проявляет \textit{вязкопластические свойства} (т.е. начинает течь лишь тогда, когда интенсивность касательных напряжений --- функция от второго инварианта девиатора тензора напряжений --- достигает фиксированного ненулевого значения). Например, жидкость Бингама--Шведова,
которая при простом сдвиге описывается соотношениями
\[
e_{xy}=\begin{cases}
\dfrac{1}{2\mu}(|\tau_{xy}|-\tau_0)\cdot\dfrac{\tau_{xy}}{|\tau_{xy}|}&\text{при }|\tau_{xy}|\geqslant\tau_0,\\[6pt]
0&\text{при }|\tau_{xy}|<\tau_0,
\end{cases}
\]
начинает движение в цилиндрическом капилляре радиуса $R$ при превышении градиентом давления значения $|\grad p|=2\tau_0/R$ ($\tau_0$ --- пороговое касательное напряжение)\footnote[1]{Уточним, что в движение приходит пристеночный слой жидкости; в центральной части капилляра при любых, даже очень больших градиентах некоторая часть жидкости движется как твердое тело.}. При моделировании пористой среды набором параллельных каналов с различными радиусами отсюда получается закон фильтрации, качественно схожий с \eqref{q2:eq-4-4}, причем предельный градиент определяется радиусом наибольшего капилляра.
\source{4}{23--24}
\end{textquote}

\end{document}
